\chapter{Математическая модель динамики опухолевых процессов}

Для изучения сложных биологических систем, таких как взаимодействие опухоли с
иммунной системой, широко применяется математическое моделирование
\cite{murray_mathematical_biology, perelson1997}. Оно позволяет изучать известные
биологические процессы, перенося их на плоскость математических уравнений, что
позволяет возможность проанализировать поведение моделируемой системы и
прогнозировать её отклик на внешние воздействия, не прибегая к потенциально
рискованным экспериментам с живыми организмами. В данной главе описана
математическая модель, которая описывает динамику опухолевых процессов во времени.

\section{Cтруктура математической модели и схема взаимодействия ее компонентов}

Рассматриваемая математическая модель основывается на классическом
макроскопическом подходе и основана на системе обыкновенных дифференциальных
уравнений (ОДУ). Использование аппарата дифференциальных уравнений позволяет
четко формализовать биологические процессы, заложенные в основу математической
модели, и с помощью математического аппарата анализировать поведение
рассматриваемых величин во времени.

Для описания взаимосвязей между компонентами моделируемой системы ниже
приведена концептуальная схема их взаимодействия (рисунок \ref{fig:model_scheme}).
Как показано на рисунке, структура модели включает три взаимодействующих субъекта:
опухолевые клетки ($T$), иммунные эффекторные клетки ($I$) и цитокины ($C$).
Опухолевые клетки на начальном этапе своего развития в организме выделяют
опухолевые антигены, которые распознаются иммунными клетками как патогенные и
запускают активацию иммунных клеток. Это запускает процесс уничтожения
опухолевых клеток иммунными эффекторными клетками. В свою очередь цитокины
служат сигнальными белковыми молекулами, сообщающими эффекторным клеткам о
наличии опухоли в организме и необходимости мобилизации иммунного ответа для
борьбы с этой опухолью.

\begin{figure}[ht!]
  \centering
  \includesvg[width=0.6\textwidth]{fig/schema.svg}
  \caption{Концептуальная графическая схема взаимодействия
  компонентов математической модели.}
  \label{fig:model_scheme}
\end{figure}

На представленной выше графической схеме (см. рисунок \ref{fig:model_scheme})
наглядно изображены биологические процессы, описанные следующей системой уравнений
(\ref{eq:ode_system}):

\begin{equation}
  \label{eq:ode_system}
  \begin{cases}
    \frac{dT}{dt} = aT (1 - bT) - (c + \eta_c(t)) TI, \\
    \frac{dI}{dt} = (d + pC(t)) TI - (\mu - \eta_{\mu}(t))I + s_A(t), \\
    \frac{dC}{dt} = s_C(t) - \lambda C.
  \end{cases}
\end{equation}
Динамика популяции опухолевых клеток в моделируемой системе описывается первым уравнением.
Слагаемое $aT(1-bT)$ выражает рост опухолевых клеток по логистическому закону
\cite{logistic_growth}.
Параметр $a$ показывает, как быстро клетки делятся, а параметр $b$ ограничивает этот рост,
поскольку с увеличением популяции достигается максимальная ёмкость среды $K=1/b$.
Слагаемое $-(c+\eta_c(t))TI$ отражает уничтожение опухолевых клеток иммунной системой.
Здесь $c$ -- скорость уничтожения клеток, а функция $\eta_c(t)$ показывает, как терапия
может её усилить.

Второе уравнение описывает динамику иммунной эффекторной популяции во времени.
Слагаемое $(d+pC(t))TI$ описывает прирост эффекторных клеток -- их активацию и
последующее деление за счет прямого столкновения с антигенами опухоли.
Интенсивность этого отклика определяется константой скорости активации иммунных
клеток и присутствием цитокинов $C(t)$, резко повышающих чувствительность
эффекторных клеток к чужеродным опухолевым клеткам с коэффициентом $p$,
обозначающим константу усиления пролиферации иммунных эффекторных клеток под
действием цитокинов.
Параметр $\mu$ обозначает скорость гибели эффекторных клеток за счет их
естественной гибели в процессе естественного старения и физиологического отмирания.
При проведении терапии этот показатель может быть увеличен ввиду проведения
внешнего терапевтического воздействия, направленного на снижение скорости
гибели иммунных клеток (параметр $\eta_{\mu}(t)$). Также для учета проведения
адоптивной иммунотерапии в рассматриваемой модели введен параметр $s_A(t)$,
обозначающий скорость адоптивного переноса иммунных клеток.

Третье уравнение описывает динамику цитокиновой популяции ($C$) во времени.
Скорость её изменения определяется константой $\lambda$, которая обозначает
темп распада цитокинов в процессе их естественного старения
\cite{cytokine_pharmacokinetics}. А параметр
$s_C(t)$ отвечает за интенсивность введения новых цитокинов в организм в рамках
соответствующей терапии цитокинами. Это позволяет смоделировать и оценить
эффективность различных режимов иммунотерапии цитокинами \cite{il2_pharmacokinetics}.
Однако стимуляция цитокинами в рамках соответствующей терапии сопряжена с
серьезными ограничениями \cite{cytokines_in_cancer_sys}. Из-за их высокой скорости
распада (параметр $\lambda$) для поддержания терапевтической концентрации
требуются значительные дозы препарата цитокина, что повышает общую токсическую
нагрузку на организм.

Совместное моделирование динамики опухолевой, иммунной и цитокиновой популяций
позволяет анализировать целостную картину развития опухоли и выявить
эффективные сценарии терапии, которые приведут к снижению нагрузки на организм
пациента или к полному его выздоровлению.

\section{Начальные значения параметров математической модели}

Определение начальных значений параметров математической модели необходимо для
корректной постановки начальных условий системы и её последующего численного решения.
В качестве начальных значений параметров были использованы биологически
обоснованные величины из работы \textit{Кузнецова В.А.} \cite{Kuznetsov1994}.
Такой выбор обусловлен высокой степенью структурного подобия математического
описания динамики опухолевых процессов в приведенной работе с математической
моделью, рассматриваемой в данной работе. Эти значения послужили основой для
последующих вычислительных экспериментов и помогают правильно описать
взаимосвязи между опухолевыми и иммунными клетками в моделируемой системе.
Найденные параметры системы приведены в таблице \ref{tab:base_params}.

\begin{table}[ht!]
  \centering
  \caption{Начальные значения параметров математической модели.}
  \label{tab:base_params}
  \begin{tabular}{|c|c|p{7cm}|}
    \hline
    Параметр & Значение & Описание и источник \\
    \hline
    $a$ & 0.18 день$^{-1}$ & Скорость пролиферации опухолевых клеток
    \cite{Kuznetsov1994} \\
    \hline
    $b$ & $2.0 \times 10^{-9}$ (клеток)$^{-1}$ & Параметр, обратный емкости среды
    \cite{Kuznetsov1994} \\
    \hline
    $c$ & $1 \times 10^{-7}$ мл/(клетка·день) & Константа скорости уничтожения
    опухолевых клеток (оценочное значение) \\
    \hline
    $d$ &  $1 \times 10^{-7}$ мл/(клетка·день) & Константа скорости активации
    иммунных клеток (оценочное значение) \\
    \hline
    $\mu$ & 0.0412 день$^{-1}$ & Константа скорости гибели иммунных клеток
    \cite{Kuznetsov1994} \\
    \hline
    $p$ & $4.0 \times 10^{-8}$ мл/(клетка·день) & Усиление пролиферации под
    действием цитокинов \\
    \hline
    $\lambda$ & 20.0 день$^{-1}$ & Константа скорости распада цитокинов
    \cite{cytokine_decay_rate} \\
    \hline
    $T_0$ & $5 \times 10^{5}$ кл/мл & Начальная концентрация опухолевых клеток
    \cite{Kuznetsov1994} \\
    \hline
    $I_0$ & $3.2 \times 10^{5}$ кл/мл & Начальная концентрация иммунных клеток
    \cite{Kuznetsov1994} \\
    \hline
    $C_0$ & 0.1 нг/мл & Начальная концентрация цитокинов \\
    \hline
  \end{tabular}
\end{table}

Величина скорости распада цитокинов взята равной $\lambda = 20.0$ день$^{-1}$,
что основано на известных фармакокинетических свойствах интерлейкина-2
\cite{il2_pharmacokinetics} -- ключевого белка, который контролирует
рост и активацию T-лимфоцитов, составляющих иммунные клетки. Такое
большое значение отражает короткий период распада этого белка в организме. Коэффициент
усиления пролиферации под действием цитокинов $p$ выбран равным $4.0
\times 10^{-8}$ мл/(кл$\cdot$сут) -- этот параметр подобран из диапазона
$[0,10d]$ для сохранения биологической стабильности уравнений.
Значение $C_0 = 0.1$ нг/мл описывает фоновое присутствие сигнальных молекул цитокина
в тканях до активации опухолевого процесса \cite{initial_cytokines}.

Как показано в таблице \ref{tab:base_params}, большинство параметров
взято из работы \textit{Кузнецова В.А.} \cite{Kuznetsov1994}, которая
признана и широко цитируется в области макроскопического
моделирования онкоиммунологии. Особое внимание стоит обратить
на параметры $c$ и $d$, так как именно они формируют баланс сил в
системе взаимодействия опухоли и иммунитета. В реальных
клинических условиях эти значения могут сильно изменяться -- иногда на несколько порядков,
что делает важным применение алгоритмов для поиска оптимальных
значений параметров модели, которые будут рассмотрены далее. Начальные значения
$T_0$ и $I_0$ тоже соответствуют типичным данным из
экспериментальных моделей \textit{in vivo}.

\section{Выбор численного метода решения системы дифференциальных уравнений и
его программная реализация}

Для численного решения систем нелинейных ОДУ, которые описывают сложные биологические
процессы \cite{nonlinear_dynamics}, важно выбрать алгоритм,
способный обеспечить высокую точность при рациональном расходе вычислительных ресурсов.
В практической части работы для моделирования динамики взаимодействия
опухолевых и иммунных клеток использовался Дормана-Принса, известный также как метод
Рунге-Кутты 4(5) порядка \cite{yang_runge_kutta}).

Прежде чем переходить к описанию итерационных процедур, рассмотрим общую форму
решаемой системы ОДУ первого порядка (\ref{eq:general_ode}):
\begin{equation}
  \label{eq:general_ode}
  \frac{dy}{dt} = f(t, y),
\end{equation}
в которой $y$ означает вектор зависимых переменных, включающий концентрации
популяций $[T, I, C]$, переменная $t$ обозначает время, а вектор-функция $f(t,
y)$ определяет правые части системы, заданные биологическими законами
пролиферации и гибели клеток.

В отличие от классических схем с фиксированным шагом, используемый в работе
метод Дормана-Принса опирается на принцип вложенных формул Рунге-Кутты, что позволяет
оценивать локальную погрешность вычислений непосредственно в процессе
интегрирования. Для нахождения следующего значения вектора состояния $y_{n+1}$
в методе Дормана-Принса используется формула (\ref{eq:rk_update}):
\begin{equation}
  \label{eq:rk_update}
  y_{n+1} = y_n + h \sum_{i=1}^{7} b_i k_i.
\end{equation}

Для определения временного шага алгоритма на каждом этапе рассчитываются
коэффициенты наклона $k_i$, описываемые системой уравнений (\ref{eq:rk_stages}):
\begin{equation}
  \label{eq:rk_stages}
  \begin{cases}
    k_1 &= f(t_n, y_n), \\
    k_2 &= f(t_n + c_2 h, y_n + h(a_{21} k_1)), \\
    k_3 &= f(t_n + c_3 h, y_n + h(a_{31} k_1 + a_{32} k_2)), \\
    \dots \\
    k_7 &= f(t_n + c_7 h, y_n + h(a_{71} k_1 + \dots + a_{76} k_6)).
  \end{cases}
\end{equation}

Ключевым преимуществом метода Дормана-Принса является автоматический выбор
величины шага по временной сетке $h$. Процесс интегрирования начинается с
задания начальных условий $y_0 = [T_0, I_0, C_0]$ и продолжается до достижения
конечного момента времени $t_{end}$. При превышении вычисленного порога
разности в решениях системы с использованием четвертого
и пятого порядков точности численного метода, шаг алгоритма автоматически
уменьшается, что обеспечивает сохранение заданного уровня точности. Если же
погрешность находится в допустимых пределах, шаг увеличивается -- это позволяет
существенно ускорить расчеты на интервалах временной сетки, где переменные
изменяются с более медленной скоростью. Такой подход позволяет минимизировать
погрешность, получаемую в процессе численного решения, что позволяет
производить более точное моделирование рассматривомой системы.
